\documentclass[11pt]{article}
\usepackage[margin=1in]{geometry}
\usepackage[T1]{fontenc}
\usepackage{lmodern,microtype,hyperref,enumitem,amsmath,amssymb,xcolor,booktabs}
\usepackage[numbers,sort&compress]{natbib}
\hypersetup{colorlinks=true,urlcolor=blue,citecolor=blue,linkcolor=blue}
\title{Q30: Inside the Framework: the PyTorch Compiler and Autodiff Memory}
\author{DLS 2026 --- module author: Kukla}
\date{October 2026}
\usepackage{url}
\begin{document}
\maketitle

\section*{Learning objectives}
\begin{itemize}[leftmargin=1.4em]
\item Explain the layers of the \texttt{torch.compile} stack and what each does.
\item Contrast tape-based (reverse-mode) and functional autodiff and their memory behavior.
\item Reason about the activation-memory cost of training and how rematerialization trades it for compute.
\end{itemize}

\begin{quote}\small\textbf{Companion videos (watch in order):}\begin{itemize}\setlength\itemsep{0pt}\item \emph{What is Automatic Differentiation?} --- Ari Seff (14:24). \url{https://www.youtube.com/watch?v=wG_nF1awSSY}\item \emph{torch.compile: The Missing Manual} --- PyTorch (15:15). \url{https://www.youtube.com/watch?v=rew5CSUaIXg}\end{itemize}\end{quote}

\section{The question}
Student response \#29 asks how the PyTorch compiler works and what its layers are, and why functional
autodiff can use far more memory than tape-based autodiff. These are two of the most ``internals'' questions
in the whole survey, and both reward a precise, mechanism-level answer.

\section{Eager vs.\ compiled execution}
By default PyTorch runs \textbf{eagerly}: each operation dispatches and executes immediately, which is easy
to debug but leaves performance on the table (Python overhead, no cross-op fusion). \texttt{torch.compile}
(introduced in PyTorch 2.0) \textbf{just-in-time compiles} the model to recover that performance while
keeping eager semantics \citep{ansel2024pytorch2}. It does this in a stack of distinct layers.

\section{The torch.compile stack, layer by layer}
\begin{itemize}[leftmargin=1.4em]
\item \textbf{TorchDynamo} hooks CPython's frame evaluation to \emph{trace} Python bytecode into an FX graph.
When it hits something it cannot trace (a data-dependent branch, an unsupported call) it inserts a
\textbf{graph break}, running that part eagerly and resuming tracing after---so compilation degrades
gracefully rather than failing.
\item \textbf{AOTAutograd} traces \emph{both} the forward and the backward graph ahead of time, so the
compiler can optimize training, not just inference, and can decide jointly what to save vs.\ recompute.
\item \textbf{A backend compiler}, by default \textbf{TorchInductor}, lowers the graph to fast kernels:
it performs \textbf{operator fusion} (combining elementwise ops to cut memory traffic) and generates
\textbf{Triton} kernels for the GPU (and C++/OpenMP for CPU).
\item \textbf{Triton} is the GPU kernel language Inductor emits---a Python-like DSL that compiles to
efficient GPU code without hand-writing CUDA.
\end{itemize}
The design philosophy is ``capture what you can, fall back where you must,'' which is why
\texttt{torch.compile} works on real models with messy Python.

\section{Two flavors of autodiff}
Automatic differentiation comes in two main styles. \textbf{Reverse-mode / tape-based} autodiff (classic
PyTorch \texttt{autograd}) runs the forward pass while recording a \emph{tape} of operations and their saved
inputs, then walks the tape backward to accumulate gradients. \textbf{Functional} autodiff (JAX-style,
\texttt{functorch}/\texttt{torch.func}) instead transforms a pure function into a new function that computes
gradients, composing transformations (\texttt{grad}, \texttt{vmap}, \texttt{jit}) functionally.

\section{Why functional AD can use more memory}
The crux of the student's question. Reverse-mode AD must keep the intermediate \textbf{activations} needed
by the backward pass alive from the forward pass---this activation memory, not the parameters, usually
dominates training memory, and it grows with network depth and batch size. A naive \emph{functional}
formulation can be \emph{worse}: without the careful liveness analysis and in-place bookkeeping that a mature
tape implementation does, composed transforms may materialize and retain more intermediate tensors (e.g.\
holding entire intermediate structures a transform could in principle free), inflating peak memory. The
practical fix in both worlds is \textbf{gradient checkpointing / rematerialization}: save only a subset of
activations and \emph{recompute} the rest during the backward pass, trading extra compute for a large memory
reduction---often turning an out-of-memory model into a trainable one at the cost of a second partial
forward. The companion code (\texttt{code/autodiff\_memory.py}) measures saved-activation count with and
without checkpointing to make the trade concrete.

\section{Why compile at all: a worked example of fusion}
To see why the compiler stack earns its complexity, consider a trivial sequence: \texttt{y = relu(x @ W +
b)}. In eager mode this is at least three separate GPU kernel launches (matmul, bias-add, relu), and the
intermediate results are written to and read back from GPU memory between each---memory traffic, not
arithmetic, is often the bottleneck. TorchInductor can \textbf{fuse} the bias-add and relu (and sometimes
more) into a single generated Triton kernel, so the intermediate never leaves fast on-chip memory and the
Python dispatch overhead disappears. Multiply that saving across every layer of a deep network and the
speedup is substantial---frequently 1.3--2$\times$ on real training, more on inference---for a one-line
change (\texttt{model = torch.compile(model)}). The graph-break mechanism is what makes this safe on real
code: when Dynamo meets a construct it cannot trace, it compiles the traceable regions and runs the rest
eagerly, so you get most of the speedup without having to rewrite your model to be ``compiler-friendly.''

\section{The memory picture in one sentence}
Tie the two halves together: the compiler (via AOTAutograd) sees the forward and backward graphs jointly, so
it can reason about exactly which activations are needed for the backward pass and even apply
rematerialization automatically---meaning the ``why does my training run out of memory'' question (dominated
by saved activations) and the ``how does the compiler help'' question are, at the deepest level, the same
question about the forward--backward activation lifetime.

\section{Summary}
\texttt{torch.compile} is a layered JIT: Dynamo traces bytecode (with graph breaks), AOTAutograd captures
forward+backward, Inductor fuses and emits Triton/C++ kernels. Autodiff is reverse-mode/tape-based (records
a tape, walks it back) or functional (transforms pure functions). Activation memory dominates training; a
naive functional formulation can retain more intermediates than a tuned tape, and checkpointing trades
recompute for memory to tame it.

\bibliographystyle{plainnat}
\bibliography{bibliography}
\end{document}
